% =============================================================================
%  Résumé (français) et Abstract (anglais)
%
%  Le résumé en arabe est dans frontmatter/resume_ar.tex : il exige un moteur
%  XeLaTeX ou LuaLaTeX (voir les instructions en tête de ce fichier). Ne
%  l'inclure que si le projet compile avec l'un des deux.
% =============================================================================

\chapter*{Résumé}
\addcontentsline{toc}{chapter}{Résumé}
\markboth{RÉSUMÉ}{RÉSUMÉ}

Ce projet de fin d'année, réalisé à distance pour EURAFRIC Information, porte sur
l'allocation d'un capital entre plusieurs actifs financiers~: comment répartir
une somme entre quatre valeurs de la Bourse de Casablanca et cinq fonds
indiciels internationaux de manière à obtenir le meilleur rendement possible pour
un risque maîtrisé.

La théorie classique répond à cette question depuis 1952, mais échoue en pratique
pour quatre raisons structurelles~: les quantités qu'elle estime sont bruitées,
le comportement des marchés change au cours du temps, la diversification
disparaît précisément pendant les crises, et les performances mesurées sur le
passé sont systématiquement optimistes. Le travail construit une chaîne complète
qui traite ces quatre problèmes~: préparation et validation des données en trois
couches, moteur de simulation temporelle dont l'absence de vision du futur est
vérifiée par des tests, puis quatre familles de modèles ajoutées une à une —
rétrécissement de la matrice de covariance, détection non supervisée de régimes
de marché par modèle de Markov caché, prédiction de rendement par ensembles
d'arbres, et optimisation tenant compte des coûts de transaction. L'évaluation
est traitée comme un objet d'ingénierie à part entière~: sélection strictement
antérieure (l'entraînement précède toujours la validation, purge et embargo à la
frontière), segment de test gelé, et \textbf{comparaisons pairées} par bootstrap
à blocs — le test de la \emph{différence} entre deux stratégies, que des
intervalles marginaux ne fournissent pas.

Les résultats sont rapportés selon leur force probante. Sur cinq fenêtres de
crise externes, le détecteur de régime est plus souvent dans son état « baissier »
qu'en dehors des fenêtres~; cette association est présentée comme exploratoire,
car cinq événements et une définition de l'état par le rendement ne permettent
pas une conclusion confirmatoire. Le système complet présente un écart observé de
6,2~\% de ratio de Sharpe net sur le portefeuille visé et de $-1{,}6$~\% sur
l'univers international. \textbf{Ces chiffres intègrent une exposition de change
USD/MAD non couverte} — les actifs BVC sont libellés en dirhams et les ETF en
dollars —~: il s'agit d'une exposition économique matérielle portée par tous les
résultats de portefeuille, non d'une réserve de forme. Sur l'instantané publié,
aucun des 1\,188 ajustements — quatre stratégies évaluées, 297 dates de
rééquilibrage — n'a utilisé de repli d'estimateur~: chaque résultat rapporté
provient du modèle que son étiquette désigne, mesure reproductible et versionnée
plutôt qu'hypothèse. Sous validation strictement antérieure, les huit
comparaisons pairées menées sur la fenêtre de test gelée ont toutes un intervalle
contenant zéro~: \textbf{aucune surperformance n'est établie}, ni pour les
modèles de signal contre la stratégie à régimes, ni contre l'équipondéré. Ce
n'est pas davantage une preuve d'équivalence, faute de test dédié. Cinq pistes indépendantes ont
été explorées et fermées, établissant notamment que précision de prédiction et
performance de portefeuille sont deux objectifs distincts. Enfin, le plafond de
position imposé comme règle métier s'est révélé le régulariseur le plus efficace
du système. Le livrable comprend un tableau de bord de recherche à deux pages et
une interface REST de consultation, couverts par 390~tests automatisés.

\vspace{0.6em}
\noindent\textbf{Mots-clés~:} optimisation de portefeuille, apprentissage
automatique, modèle de Markov caché, matrice de covariance, backtesting,
validation hors échantillon, Bourse de Casablanca, gestion du risque.

\chapter*{Abstract}
\addcontentsline{toc}{chapter}{Abstract}
\markboth{ABSTRACT}{ABSTRACT}

This final-year project, carried out remotely for EURAFRIC Information, addresses capital
allocation across financial assets: how to split a sum between four Casablanca
Stock Exchange equities and five international index funds so as to obtain the
best possible return for a controlled level of risk.

Classical theory has answered this question since 1952, but fails in practice for
four structural reasons: the quantities it estimates are noisy, market behaviour
changes over time, diversification vanishes precisely during crises, and
performance measured on past data is systematically optimistic. This work builds
a complete pipeline addressing all four: a three-layer data preparation and
validation stack, a walk-forward simulation engine whose inability to see the
future is enforced by tests rather than asserted, and four model families added
one at a time — covariance shrinkage, unsupervised market-regime detection via a
hidden Markov model, return prediction with tree ensembles, and
transaction-cost-aware optimisation. Evaluation is treated as an engineering
artefact in its own right: strictly forward-only selection (training always
precedes validation, with a purge and an embargo at the boundary), a frozen test
segment, and \textbf{paired} block-bootstrap comparisons — a test of the
\emph{difference} between two strategies, which marginal intervals do not provide.

Results are reported by strength of evidence. Across five external crisis
windows, the regime detector is more often in its ``bear'' state than outside
those windows; this association is reported as exploratory because five events
and a return-based state definition cannot support a confirmatory conclusion.
The full system has an observed 6.2\% net-Sharpe difference on the target
portfolio and a $-1.6$\% difference on the international universe.
\textbf{These figures embed an unhedged USD/MAD currency exposure} — BVC assets
are MAD-denominated and the ETFs USD — a material economic exposure carried by
every portfolio result, not a formality. On the released snapshot none of the
1\,188 fits — four evaluated strategies, 297 rebalance dates — used an estimator
fallback: every reported result was produced by the model its label names, a
reproducible versioned measurement rather than an assumption. Published
under forward-only selection, all eight paired comparisons on the frozen test
window have intervals containing zero: \textbf{no outperformance is established},
for either signal model against either the regime strategy or equal weighting.
Nor is this evidence of equivalence, absent a dedicated test.
Five independent avenues for improvement were explored and closed, establishing
among other things that prediction accuracy and portfolio performance are
distinct objectives. Finally, the position cap imposed as a business rule turned
out to be the system's most effective regulariser. Deliverables include a
two-page research dashboard and a read-only REST interface, covered by 390
automated tests.

\vspace{0.6em}
\noindent\textbf{Keywords:} portfolio optimisation, machine learning, hidden
Markov model, covariance estimation, backtesting, out-of-sample validation,
Casablanca Stock Exchange, risk management.
